% --- 题目 ---
\problem[P229-2 (22-1,3)]{设 $\displaystyle X_1, X_2, \dots, X_n$ 为来自均值为 $\displaystyle \theta$ 的指数分布总体的简单随机样本, $\displaystyle Y_1, Y_2, \dots, Y_m$ 为来自均值为 $\displaystyle 2\theta$ 的指数分布总体的简单随机样本, 且两样本相互独立, 其中 $\displaystyle \theta (\theta > 0)$ 为未知参数. 利用样本 $\displaystyle X_1, X_2, \dots, X_n, Y_1, Y_2, \dots, Y_m$, 求 $\displaystyle \theta$ 的最大似然估计量 $\displaystyle \hat{\theta}$, 并求 $\displaystyle D(\hat{\theta})$.
}
\ansat{371；【十年真题】 - 考点一：矩估计与最大似然估计 - 2}

% --- 题目 ---
\problem[P229-3 (20-1,3)]{设某种元件的使用寿命 $\displaystyle T$ 的分布函数为
\[ F(t) = \begin{cases} 1 - \mathrm{e}^{-(\frac{t}{\theta})^m}, & t \ge 0, \\ 0, & \text{其他}, \end{cases} \]
其中 $\displaystyle \theta, m$ 为参数且大于零.
\begin{enumerate}
\item[(1)] 求概率 $\displaystyle P\{T>t\}$ 与 $\displaystyle P\{T>t+s | T>s\}$, 其中 $\displaystyle s>0, t>0$;
\item[(2)] 任取 $\displaystyle n$ 个这种元件做寿命试验, 测得它们的寿命分别为 $\displaystyle t_1, t_2, \dots, t_n$. 若 $\displaystyle m$ 已知, 求 $\displaystyle \theta$ 的最大似然估计值 $\displaystyle \hat{\theta}$.
\end{enumerate}}
\ansat{371；【十年真题】 - 考点一：矩估计与最大似然估计 - 3}

% --- 题目 ---
\problem[P229-4 (19-1,3)]{设总体 $\displaystyle X$ 的概率密度为
\[ f(x; \sigma^2) = \begin{cases} \displaystyle \frac{A}{\sigma} \mathrm{e}^{-\frac{(x-\mu)^2}{2\sigma^2}}, & x \ge \mu, \\ 0, & x < \mu, \end{cases} \]
其中 $\displaystyle \mu$ 是已知参数, $\displaystyle \sigma > 0$ 是未知参数, $\displaystyle A$ 是常数, $\displaystyle X_1, X_2, \dots, X_n$ 是来自总体 $\displaystyle X$ 的简单随机样本.
\begin{enumerate}
\item[(1)] 求 $\displaystyle A$;
\item[(2)] 求 $\displaystyle \sigma^2$ 的最大似然估计量.
\end{enumerate}}
\ansat{372；【十年真题】 - 考点一：矩估计与最大似然估计 - 4}
\vspace{8em}

% --- 题目 ---
\problem[P229-5 (18-1,3)]{设总体 $\displaystyle X$ 的概率密度为
\[ f(x; \sigma) = \frac{1}{2\sigma} \mathrm{e}^{-\frac{|x|}{\sigma}}, \quad -\infty < x < +\infty, \]
其中 $\displaystyle \sigma \in (0, +\infty)$ 为未知参数, $\displaystyle X_1, X_2, \dots, X_n$ 为来自总体 $\displaystyle X$ 的简单随机样本. 记 $\displaystyle \sigma$ 的最大似然估计量为 $\displaystyle \hat{\sigma}$.
\begin{enumerate}
\item[(1)] 求 $\displaystyle \hat{\sigma}$;
\item[(2)] 求 $\displaystyle E\hat{\sigma}$ 和 $\displaystyle D\hat{\sigma}$.
\end{enumerate}}
\ansat{372；【十年真题】 - 考点一：矩估计与最大似然估计 - 3}
